\begin{theorem}[Vacancy in an actual horizontal two-plaquette block]
\label{thm:peps_actual_horizontal_vacancy}
\lean{TNLean.PEPS.regularWalkHolonomy_rightPlaquette_eq_one_iff_treeGauge,
TNLean.PEPS.regularRegionTreeCycleResidual_middle_eq_one_iff_rightPlaquette,
TNLean.PEPS.translatedTwoPlaquetteIso_vertex_coordinates}
\leanok
\uses{thm:peps_tree_gauge_transport}
On a finite torus of width at least four and height at least three, take
the actual translated three-by-two block with its derived spanning path.
For arbitrary ordered bond operators, the right plaquette is vacant exactly
when the middle upward bond has trivial residual in the root-normalized
tree gauge. Both periodic seams are included.
This is an auxiliary finite-torus criterion for SCP10, Theorem~6.16,
lines 2271--2305; the unrestricted geometry and the complete prescribed
braid remain separate \cite{gap:rmp_peps_quantum_double_g_isometry}.
\end{theorem}
\begin{proof}\leanok
The other three sides of the right plaquette belong to the spanning path.
Apply the three-tree-edge criterion and identify the fourth edge with
the derived middle chord. Directed transport retains the sorting convention
at a periodic seam.
\end{proof}

\begin{theorem}[A fixed physical operation on arbitrary actual patch inputs]
\label{thm:peps_actual_horizontal_flux_operation}
\lean{TNLean.PEPS.IsGIsometric.exists_unitary_torusActualHorizontalFluxMove,
TNLean.PEPS.torusActualHorizontalFluxMove,
TNLean.PEPS.torusActualHorizontalFluxMove_eq_of_ne,
TNLean.PEPS.torusActualHorizontalFluxMove_middle_of_vacancy,
TNLean.PEPS.torusActualHorizontalFluxMove_up_transport_of_vacancy}
\leanok
\uses{thm:peps_actual_cycle_physical_permutation,thm:peps_actual_horizontal_vacancy}
Assume the regular virtual action and original $G$-isometric site tensors.
For each actual translated horizontal block there is one six-spin unitary,
chosen before every actual bond assignment and every outgoing boundary
configuration. It multiplies the derived middle cycle residual on the left
by the derived left residual and reconstructs with the original tree gauge.
It acts on both the actual open contraction and the actual closed state,
retaining every ordered bond other than the middle chord, including all
crossing and exterior bonds. If the right plaquette is vacant, its new
middle ordered coefficient is $k_{\mathrm{head}}\omega_{\mathrm{left}}
 k_{\mathrm{tail}}^{-1}$. Equivalently, its directed upward transport becomes
$T L B^{-1}$, where $T$, $L$, and $B$ are the original top-left rightward,
left upward, and bottom-left rightward transports, respectively.
The finite-torus scope is the same as above. This auxiliary operation uses
actual derived coordinates; it does not assert the complete prescribed braid
\cite{gap:rmp_peps_quantum_double_g_isometry}.
\end{theorem}
\begin{proof}\leanok
Apply the original-spin cycle permutation to the tree gauge and residuals
derived from each input. Reconstruction recovers the input literally.
Controlled multiplication changes only the selected middle residual.
Vacancy makes its initial residual equal to the identity. The two upward
chords have the same sorting comparison, including at the vertical seam.
The two adjacent tree-edge gradients then give the literal formula $T L B^{-1}$.
\end{proof}
